\chapter{Preliminaries}\label{chapter:preliminaries}

We begin by introducing the potential outcomes framework~\citep{rubin1974} that we will formulate 
our mediational effects in terms of. We also introduce in more detail the running example 
from~\cite{morenobetancur2018} introduced in Section~\ref{chapter:intro} that we will use to 
demonstrate the different types of effects we are interested in estimating. We first review the 
potential outcomes framework for ``simple''
causal inference problems with no mediation, where the goal is to infer the causal 
effect of a particular change in a treatment variable on an outcome variable. 

Suppose we are weighing whether a certain anti-depressant that has been developed is effective in treating
depression. To evaluate this question, we have access to observational data from a study on $N$
patients diagnosed with clinical depression indexed by $i = 1, 2, ..., N$. It is imagined that
these patients were a random sample taken from a larger population of clinically despressed patients. Each patient,
at the start of the study period, was given the option of beginning uptake of the anti-depressant. 
We have the following data:
\begin{itemize}[noitemsep,topsep=0pt]
\item The treatment variable $T_{i}$ denotes a binary treatment variable that equals $1$ if the 
unit takes a certain dose of the anti-depressant and equals $0$ otherwise. 
\item The outcome variable $Y_{i}$ represents the patient's score on a questionnaire evaluating the severity 
of depression in the patient that is filled out months after the patient begins or does not begin treatment. 
\item Covariates $C_{i}$ are unit-specific characteristics (e.g. sex, medical history) that are potentially 
related to the patient's propensity to take treatment and the strength of their depression 6 months later. 
\end{itemize}

We suppose $T_{i}$, $C_{i}$, and $Y_{i}$ are independent and identically distributed: they are drawn from 
population random variables $T$, $C$, and $Y$. Let $\mathcal{T}$, $\mathcal{Y}$, and $\mathcal{C}$ 
denote the supports of $T$, $Y$, and $C$ respectively. 

Under the potential outcomes framework, it is imagined that each unit in our sample also has
two potential values of $Y_{i}$ --- their realized outcome under the treatment level they actually received
and the outcome they would have experienced had they received the other treatment level. Let $Y_{i}(1)$
and $Y_{i}(0)$ denote the outcome values unit $i$ would obtain under each treatment level, drawn from 
$Y(1)$ and $Y(0)$ respectively. One of these
potential outcomes will always be observed and will be equal to the realized $Y_{i}$, while the other
will always be missing. 

The \textit{unit level treatment effect} of $T$ on $Y$ is equal to $Y_{i}(1) - Y_{i}(0)$ and represents 
the effect of taking anti-depressants on unit $i$'s severity of depression. This
effect is never identified for individual units because one potential outcome for each unit is always
unobserved ---~\cite{rubin1974} and~\cite{holland1986} call this ``the fundamental problem of causal 
inference''. However, different population-level causal estimands are identified under certain assumptions. 
One such estimand we focus on is the \textit{average treatment effect} (ATE), given by 
$\mathbb{E}[Y(a^{*}) - Y(a)]$.

The ATE represents the average causal effect of the treatment on the outcome in the larger population. In
our example, this represents the average difference of patients' severity of depression in worlds where they 
are induced to take the anti-depressant medication compared to worlds in which they 
do not take the medication, holding all other relevant factors constant. The ATE, and other 
important estimands, can be identified given the following assumptions~\citep{imbens2015}: 

\begin{assumption}[Consistency]\label{assm:Consistency}
$Y_{i}(a) = Y_{i}$ if unit $i$ receives treatment level $a \in \mathcal{T}$. 
\end{assumption}

\begin{assumption}[No interference]\label{assm:no_interference}
$Y_{i}(a)$ depends only on the treatment level $a$ for unit $i$; in particular, it does not depend on 
the treatment level for any other unit $j \neq i$. 
\end{assumption}

\begin{assumption}[Conditional Ignorability]\label{assm:ignorability} After controlling for some set of 
observable covariates $C_{i}$, treatment level is independent of unit potential 
outcomes: $Y_{i}(a) \indep T_{i} | C_{i} = c$ for all $a \in \mathcal{T}$, where 
$0 < \mathbb{P}[T = a | C] < 1$ for all $a \in \mathcal{T}$. 
\end{assumption}

The first two assumptions are sometimes presented together as the Stable Unit Treatment Value Assumption 
(SUTVA)~\citep{rubin1980a} and give technical conditions under which
the ATE can be identified. It requires that there is no interference between units, meaning that a unit's 
treatment status does not affect the potential outcomes of other units and vice versa. In this example, a violation
might occur if we believed depression within a household was ``contagious'', meaning that whether the people a patient
was living with received treatment or not (if they were included in the study) affected the patient's outcome.  
Additionally, there are to be no ``hidden versions'' of treatments, meaning that a given treatment 
level cannot map to multiple different treatments for a unit. For example, in the birth control example from 
Section~\ref{chapter:intro}, a violation would occur if some women who are recording as receiving the treatment take
expired birth control pills that have lower efficacy. Violations can also occur in experimental settings where differences
in the mode of treatment assignment itself --- for example, the attitude of the doctor or lab assistant administering 
treatment --- might affect potential outcomes.

% Think about cutting/moving
The extension of the SUTVA assumption to mediation analysis is not to be made lightly. In particular,~\cite{imaietal2013} 
discusses the difficulty of satisfying the consistency assumption in hypothetical 
experimental designs for identifying natural direct and indirect effects. However, violations of SUTVA are not 
the focus of this dissertation, and SUTVA and its extensions will be assumed without further comment 
in the rest of this dissertation. 

The conditional ignorability assumption~\citep{rubin1978} allows us to identify estimands like the ATE because while both 
potential outcomes for a single unit can never be observed, ignorability allows us to treat these missing potential 
outcomes as if they were ``missing at random'' because there will be no bias driving which potential outcomes are revealed
versus hidden~\citep{imbens2015} after controlling for confounders $C$. This assumption
rules out selection bias: for example, it rules out the possibility that patients with an unobserved lower desire to 
make positive lifestyle changes are also less likely to take the anti-depressants. Alternatively, a randomized experiment
where all study participants were randomly assigned to either take the anti-depressant or a placebo would also satisfy 
this assumption. Thus, the difference in a covariate-weighted average of the realized potential outcomes at both
treatment levels is an unbiased estimator of the ATE.

Throughout this paper, we will also present directed acyclic graph (DAG) representations of the causal structures 
we are developing within the potential outcomes framework. The non-parametric structural equation 
modeling or DAG framework of~\cite{pearl2000} is an alternative method 
for approaching causal inference problems that has been adapted by many applied researchers, especially in 
epidemiology~\citep{hernan2024}. It is particularly helpful for visually encoding the types of identification 
assumptions made above. Much of the early work on mediation analysis (e.g.~\citep{pearl2001,avinetal2005}) was
conducted within the DAG framework. We find DAG representations to be helpful in gaining intuition for the types of 
complicated causal structures with many variables that are encountered in mediation analysis, but these graphs are 
presented as illustrative aids only. 

\begin{figure}[ht]
\centering
\begin{tikzpicture}

  % Define nodes
  \node[obs] (A) {$\mathbf{T}$};
  \node[obs, right=1.5cm of A]  (Y) {$\mathbf{Y}$};
  \node[obs, left=1.5cm of A]  (C) {$\mathbf{C}$};

  % Connect the nodes
  \edge {A} {Y} ; %
  \draw[->] (C) to[out=315, in=225] (Y); %
  \edge {C} {A} ; %

\end{tikzpicture}
\caption{DAG encoding the conditional ignorability assumption}
\label{fig:ignorability_dag}
\end{figure}

% Potentially add discussion of the difference between randomized experiments and observational studies here; can
% be used to emphasize the difficulty of justifying the assumptions for natural effect identification
